% from ega4-1-fr.tex
\oldpage[IV-1]{238}

\begin{env}[1.8.1]
\label{IV.1.8.1-fr}
Soit $X$ un préschéma~; on sait que tout ouvert quasi-compact dans $X$ est
réunion finie d'ouverts affines, et réciproquement. Pour qu'un ouvert $U$ de
$X$ soit \emph{rétrocompact} dans $X$ (0\textsubscript{III}, 9.1.1), il faut
et il suffit donc que pour tout \emph{ouvert affine} $V$ de $X$, $U\cap V$
soit quasi-compact. Lorsque $X$ est \emph{quasi-séparé} (1.2.1), tout ouvert
quasi-compact dans $X$ est rétrocompact dans $X$ (1.2.7), donc toute partie
localement constructible de $X$ est rétrocompacte dans $X$
(0\textsubscript{III}, 9.1.13)~; si \emph{de plus} $X$ est
\emph{quasi-compact}, il y a identité entre parties constructibles et parties
localement constructibles dans $X$ (0\textsubscript{III}, 9.1.12), et entre
parties ouvertes constructibles et parties ouvertes quasi-compactes
(0\textsubscript{III}, 9.1.5).
\end{env}

\begin{proposition}[1.8.2]
\label{IV.1.8.2-fr}
Soient $X$, $Y$ deux préschémas, $f:X\to Y$ un morphisme. Pour toute partie
constructible (resp. localement constructible) $Z$ de $Y$, $f^{-1}(Z)$ est
une partie constructible (resp. localement constructible) de $X$.
\end{proposition}

\begin{proof}
Supposons $Z$ localement constructible~; pour tout $x\in X$, il y aura un
voisinage ouvert affine $V$ de $f(x)$ tel que $Z\cap V$ soit constructible
dans $V$~; si la proposition est démontrée pour les parties constructibles,
$f^{-1}(Z)\cap f^{-1}(V)$ sera constructible dans $f^{-1}(V)$, donc
$f^{-1}(Z)$ sera localement constructible. Il suffit par suite de considérer
le cas où $Z$ est constructible, et compte tenu de (0\textsubscript{III},
9.1.3), on est ramené au cas où $Z$ est ouvert et \emph{rétrocompact} dans
$Y$, autrement dit tel que l'injection canonique $Z\to Y$ soit un morphisme
quasi-compact~; il suffit alors de voir que $f^{-1}(Z)$ est
\emph{rétrocompact dans $X$}, autrement dit tel que l'injection canonique
$f^{-1}(Z)\to X$ soit un morphisme quasi-compact~; mais cela résulte de
(1.1.2, (iii)), puisque $f^{-1}(Z)=Z\times_YX$ (I, 4.4.1).
\end{proof}

\begin{lemma}[1.8.3]
\label{IV.1.8.3-fr}
Soient $X$ un préschéma quasi-compact et quasi-séparé, $Z$ une partie
constructible de $X$. Il existe alors un schéma affine $X'$ et un morphisme
de présentation finie $f:X'\to X$ tels que $f(X')=Z$.
\end{lemma}

\begin{proof}
Comme $X$ est quasi-compact, il est réunion finie d'ouverts affines $X_i$
($i\in I$), et comme $X$ est quasi-séparé, il résulte de (1.2.7) que chacune
des immersions canoniques $g_i:X_i\to X$ est de présentation finie~; on en
conclut que si $f_i:X'_i\to X_i$ est un morphisme de présentation finie, il
en est de même de $g_i\circ f_i:X'_i\to X$ (1.6.2), et si $X'$ est le
préschéma somme des $X'_i$, le morphisme $h:X'\to X$ qui coïncide avec
$g_i\circ f_i$ dans chaque $X'_i$ est lui aussi de présentation finie
(1.6.5). Comme $Z\cap X_i$ est constructible dans $X_i$
(0\textsubscript{III}, 9.1.8), on voit qu'on est ramené à prouver le lemme
lorsque $X_i$ est \emph{affine} d'anneau $A$. Comme $Z$ est alors réunion
finie d'ensembles de la forme $U\cap\complement V$, où $U$ et $V$ sont des
ouverts quasi-compacts (0\textsubscript{III}, 9.1.3), on voit, en considérant
encore une somme convenable de préschémas, qu'on peut se ramener au cas où
$Z=U\cap\complement V$~; comme d'ailleurs $U$ est réunion finie d'ouverts
affines, on peut même se borner au cas où $U$ est \emph{affine}. Comme $V$
est quasi-compact, il est réunion finie d'ouverts de la forme $X_{f_i}$, où
$f_i\in A$~; soit $\mathfrak J$ l'idéal de $A$ engendré par les $f_i$, et
soit $Z'$ le sous-schéma fermé affine de $X$ qu'il définit (et qui est par
construction de présentation finie sur $X$)~; on a par définition
$V=\complement Z'$, donc $U\cap\complement V=U\cap Z'$. Considérons le
schéma affine $X'=Z'\times_XU$

